% Source: https://artofproblemsolving.com/community/c3h1660973_explain_this
% Stable identity: aops-topic-1660973-q01
At a steady speed, the distance covered in $1$ hour is the total distance divided by the number of hours:
\[
\dfrac{89.1}{2.2}\text{ km}.
\]
Multiplying the numerator and the denominator of a fraction by the same nonzero number does not change its value, so multiplying both by $10$ gives
\[
\dfrac{89.1}{2.2}=\dfrac{891}{22}=40.5 .
\]
(Indeed $22\cdot40=880$ and $22\cdot0.5=11$.) The car covers $\boxed{40.5\text{ km}}$ in one hour. Check: $40.5\cdot2.2=81+8.1=89.1$.
